\documentclass[conference]{IEEEtran}
\IEEEoverridecommandlockouts
\usepackage{cite}
\usepackage{amsmath,amssymb,amsfonts}
\usepackage{algorithmic}
\usepackage{graphicx}
\usepackage{url}
\usepackage{float}
\usepackage{textcomp}
\usepackage{xcolor}
\def\BibTeX{{\rm B\kern-.05em{\sc i\kern-.025em b}\kern-.08em
    T\kern-.1667em\lower.7ex\hbox{E}\kern-.125emX}}
\begin{document}

\title{Hierarchical Sparse Routing Transformer with Cross-Sentence Consistency Gate for Fake News Detection\\}

\author{\IEEEauthorblockN{Anurag Mukherjee}
\IEEEauthorblockA{\textit{Computer Science and Engineering} \\
\textit{Jalpaiguri Govt. Engineering College}\\
Jalpaiguri, West Bengal \\
\texttt{amdgp16@gmail.com}}
\and
\IEEEauthorblockN{Subhas Barman}
\IEEEauthorblockA{\textit{Computer Science and Engineering} \\
\textit{Jalpaiguri Govt. Engineering College}\\
Jalpaiguri, West Bengal \\
\texttt{subhas.barman@gmail.com}}
\and
\IEEEauthorblockN{Rebanta Sarkar}
\IEEEauthorblockA{\textit{Computer Science and Engineering} \\
\textit{Jalpaiguri Govt. Engineering College}\\
Jalpaiguri, West Bengal \\
\texttt{sarkarrebanta.2004@gmail.com}}
\and
\IEEEauthorblockN{Arijit Roy}
\IEEEauthorblockA{\textit{Computer Science and Engineering} \\
\textit{Jalpaiguri Govt. Engineering College}\\
Jalpaiguri, West Bengal \\
\texttt{aroy6934@gmail.com}}
\and
\IEEEauthorblockN{Ardhendu Mandal}
\IEEEauthorblockA{\textit{Computer Science and Technology} \\
\textit{University of North Bengal}\\
Darjeeling, West Bengal \\
\texttt{am.csa.nbu@nbu.ac.in}}
}

\maketitle

\begin{abstract}
The proliferation of misinformation across digital news platforms poses a significant threat to public discourse and democratic processes. Existing transformer-based approaches for fake news detection predominantly rely on fine-tuned BERT variants that employ absolute positional encodings, standard feed-forward networks, and single-vector \texttt{[CLS]} classification heads, leaving several architectural inefficiencies unaddressed. In this work, we propose the \textit{FakeNews Transformer}, a modified encoder architecture incorporating eight targeted improvements: Rotary Positional Encoding (RoPE) for relative position awareness, Sparse Top-K Attention for focused claim detection, Hierarchical Sparse Routing (HSR) that assigns dense attention to lower layers and sparse attention to upper layers, a Cross-Sentence Consistency Gate (CSCG) that explicitly suppresses tokens inconsistent with the global document claim, SwiGLU feed-forward activations, a Mixture-of-Experts FFN in upper layers, Pre-Layer Normalization for training stability, and a Multi-Pool classification head combining \texttt{[CLS]}, mean-pool, and max-pool representations. Experiments on the WELFake dataset demonstrate that the proposed architecture achieves competitive performance across twelve evaluation metrics, with statistically significant improvements over a BERT baseline confirmed by McNemar's test. An ablation study quantifies the individual contribution of each modification, with CSCG and Multi-Pool yielding the largest gains. Attention rollout visualizations and gradient saliency maps further validate that the model attends to semantically meaningful tokens.
\end{abstract}

\vspace{0.5em}

\begin{IEEEkeywords}
Fake News Detection, Transformer Encoder, Rotary Positional Encoding, Sparse Attention, Mixture of Experts, Cross-Sentence Consistency, Hierarchical Routing, WELFake.
\end{IEEEkeywords}

% ─────────────────────────────────────────────────────────────────────────────
\section{Introduction}
% ─────────────────────────────────────────────────────────────────────────────

The rapid dissemination of fabricated or misleading news articles through social media and online platforms has emerged as one of the most pressing challenges in natural language processing and information security \cite{shu2017fake, zhou2020survey}. Unlike spam or phishing, which exhibit relatively stable lexical signatures, fake news is deliberately crafted to mimic the style and vocabulary of credible journalism, making surface-level pattern matching insufficient for reliable detection \cite{thorne2018fever}. The consequences of undetected misinformation range from public health crises to electoral interference, underscoring the need for accurate, interpretable, and computationally efficient automated detection systems \cite{vosoughi2018spread}.

The introduction of BERT \cite{devlin2019bert} and its derivatives has substantially advanced the state of the art in text classification, including fake news detection. These models leverage large-scale pretraining on masked language modeling objectives to acquire rich contextual representations, which are subsequently fine-tuned on downstream tasks. However, several architectural limitations constrain their effectiveness in the fake news domain. First, absolute positional encodings do not generalize well to variable-length news articles, which may range from brief headlines to multi-paragraph reports \cite{su2021rope}. Second, full self-attention distributes attention mass uniformly across all tokens, including padding and stopwords, diluting the signal from claim-bearing phrases \cite{zaheer2020bigbird}. Third, the standard \texttt{[CLS]}-only classification head discards potentially discriminative information encoded in individual token representations, such as emotionally loaded language or named entities \cite{liu2019roberta}. Fourth, uniform feed-forward networks across all layers do not exploit the hierarchical nature of linguistic processing, where lower layers capture syntactic patterns and upper layers encode semantic claims \cite{tenney2019bert}.

To address these limitations, we propose the \textit{FakeNews Transformer}, a purpose-built encoder architecture that integrates eight complementary modifications. The central architectural contributions are: (i) \textit{Hierarchical Sparse Routing} (HSR, M3), a novel strategy that assigns dense multi-head attention to the lower half of encoder layers for local syntactic processing and sparse top-k attention to the upper half for selective global semantic reasoning; (ii) a \textit{Cross-Sentence Consistency Gate} (CSCG, M4), a novel gating mechanism applied after the final encoder layer that computes a per-token consistency score with respect to the \texttt{[CLS]} representation and suppresses tokens that contradict the global document claim; and (iii) a \textit{Multi-Pool Classification Head} (M8), which concatenates the \texttt{[CLS]} token, mean-pooled, and max-pooled representations to capture complementary views of the document before classification. These three contributions are novel to this work. The remaining five modifications---RoPE (M1), Sparse Top-K Attention (M2), SwiGLU FFN (M5), Mixture-of-Experts FFN (M6), and Pre-Layer Normalization (M7)---are established techniques from the recent literature that we adapt and integrate into a unified architecture for the fake news detection task.

The model is trained and evaluated on the WELFake dataset \cite{verma2021welfake}, a large-scale benchmark comprising approximately 72,000 news articles labeled as real or fake. We report twelve evaluation metrics including accuracy, macro F1-score, Matthews Correlation Coefficient, and calibration error, along with 95\% bootstrap confidence intervals. Statistical significance of improvements over the BERT baseline is assessed using McNemar's test. An ablation study over nine model configurations isolates the contribution of each modification. Attention rollout \cite{abnar2020quantifying} and integrated gradient \cite{sundararajan2017axiomatic} visualizations provide qualitative evidence that the model attends to semantically meaningful tokens.

The remainder of this paper is organized as follows. Section~\ref{sec:related} reviews related work. Section~\ref{sec:method} describes the dataset, preprocessing pipeline, and full architecture. Section~\ref{sec:experiments} presents the experimental setup, main results, ablation study, and statistical analysis. Section~\ref{sec:analysis} provides attention visualization and error analysis. Section~\ref{sec:conclusion} concludes the paper.

% ─────────────────────────────────────────────────────────────────────────────
\section{Related Work}
\label{sec:related}
% ─────────────────────────────────────────────────────────────────────────────

\subsection{Fake News Detection}

Early approaches to automated fake news detection relied on hand-crafted linguistic features such as writing style, sentiment polarity, and readability scores, combined with classical classifiers including Support Vector Machines and Logistic Regression \cite{castillo2011credibility, potthast2018stylometric}. While these methods established useful baselines, they exhibited limited generalization across domains and were vulnerable to adversarial stylistic mimicry. Subsequent work explored knowledge-graph-based verification \cite{popat2018declare} and claim-evidence matching \cite{thorne2018fever}, which improved factual grounding but required structured external knowledge bases not always available at inference time.

The introduction of deep learning substantially advanced detection performance. Convolutional and recurrent architectures, including CNNs and LSTMs, demonstrated the ability to capture local n-gram patterns and sequential dependencies in news text \cite{wang2018eann, ma2016detecting}. Attention mechanisms further improved these models by enabling selective focus on claim-bearing phrases. However, these architectures lacked the capacity for deep bidirectional contextual reasoning that characterizes transformer-based models.

\subsection{Transformer-Based Approaches}

The release of BERT \cite{devlin2019bert} marked a turning point in fake news detection. Fine-tuned BERT models consistently outperformed prior approaches on benchmarks such as LIAR \cite{wang2017liar} and FakeNewsNet \cite{shu2020fakenewsnet}, demonstrating the value of large-scale pretraining for acquiring world knowledge and linguistic priors. Subsequent variants including RoBERTa \cite{liu2019roberta}, ALBERT \cite{lan2019albert}, and DeBERTa \cite{he2021deberta} further improved performance through improved pretraining objectives, parameter sharing, and disentangled attention. Domain-specific models such as NewsEmbed \cite{hamborg2019newsembed} and FakeBERT \cite{kaliyar2021fakebert} adapted pretraining to the news domain, yielding additional gains.

Despite these advances, standard BERT-based approaches share several architectural limitations relevant to fake news detection. Absolute positional encodings do not capture relative positional relationships between claim and evidence tokens, which is critical for detecting headline-body contradictions \cite{su2021rope}. Full self-attention over all token pairs is computationally redundant and may dilute attention from semantically critical tokens \cite{zaheer2020bigbird, beltagy2020longformer}. The \texttt{[CLS]}-only classification head may fail to capture strong local signals such as emotionally loaded words or named entities that are characteristic of fabricated content.

\subsection{Efficient and Sparse Attention}

Sparse attention mechanisms have been proposed to reduce the quadratic complexity of full self-attention and to improve focus on relevant tokens. Longformer \cite{beltagy2020longformer} combines local sliding-window attention with global attention on selected tokens. BigBird \cite{zaheer2020bigbird} extends this with random attention blocks, achieving linear complexity while preserving theoretical expressiveness. Reformer \cite{kitaev2020reformer} uses locality-sensitive hashing to approximate full attention. In contrast to these methods, which are primarily motivated by efficiency on long sequences, our Sparse Top-K Attention (M2) is motivated by the need to force the model to identify the most claim-relevant tokens in news articles, and is applied selectively to the upper encoder layers via the HSR strategy (M3).

\subsection{Positional Encoding}

Rotary Positional Encoding (RoPE) \cite{su2021rope} encodes position as a rotation in the complex plane applied to query and key vectors, such that the inner product between rotated queries and keys depends only on their relative position. This property is particularly beneficial for variable-length news articles. RoPE has been adopted in LLaMA \cite{touvron2023llama}, PaLM \cite{chowdhery2023palm}, and Mistral \cite{jiang2023mistral}, demonstrating its effectiveness across scales. We adapt RoPE for the classification setting and apply it within each attention head of our encoder.

\subsection{Feed-Forward Network Variants}

SwiGLU \cite{shazeer2020glu} replaces the standard ReLU feed-forward network with a gated variant: $\text{FFN}(x) = (xW_1 \odot \text{SiLU}(xW_3))W_2$, where the gate branch learns to suppress irrelevant activations. SwiGLU has been shown to consistently outperform ReLU and GeLU activations on language modeling benchmarks and is used in PaLM and LLaMA. Mixture-of-Experts (MoE) FFNs \cite{fedus2022switch} replace a single FFN with a set of expert networks and a learned router, enabling the model to develop specialized sub-networks for different input types. We apply soft-routing MoE in the upper encoder layers, where expert specialization is most beneficial for distinguishing semantic patterns associated with real and fake news.

\subsection{Layer Normalization Placement}

Xiong et al.\ \cite{xiong2020layer} demonstrated that Pre-Layer Normalization (Pre-LN), which applies layer normalization before each sub-layer rather than after, substantially improves training stability and convergence speed, particularly on smaller datasets and with limited warmup. This is directly relevant to our setting, where training is conducted on a single GPU without extended warmup schedules.

% ─────────────────────────────────────────────────────────────────────────────
\section{Proposed Methodology}
\label{sec:method}
% ─────────────────────────────────────────────────────────────────────────────

\subsection{Dataset}

We conduct all experiments on the WELFake dataset \cite{verma2021welfake}, a large-scale benchmark for binary fake news classification. The dataset aggregates articles from four sources---Kaggle, McIntire, PolitiFact, and GossipCop---yielding a total of 72,134 samples, of which 35,028 are labeled as real and 37,106 as fake. Each sample consists of a news title and body text. The dataset is split into training, validation, and test sets using stratified sampling to preserve the class prior:
\[
\mathcal{D} = \mathcal{D}_{\text{train}} \cup \mathcal{D}_{\text{val}} \cup \mathcal{D}_{\text{test}}, \quad |\mathcal{D}_{\text{train}}| : |\mathcal{D}_{\text{val}}| : |\mathcal{D}_{\text{test}}| = 0.8 : 0.1 : 0.1.
\]
The dataset is downloaded automatically from HuggingFace Hub on first run and cached locally.

\subsection{Preprocessing and Tokenization}

Each article is represented as a structured token sequence comprising three segments: the headline (segment~0), the body text (segment~1), and optional metadata such as source and date (segment~2). The full sequence is constructed as:
\[
x = [\texttt{[CLS]},\; \underbrace{w_1^h, \dots, w_{T_h}^h}_{\text{headline}},\; \underbrace{w_1^b, \dots, w_{T_b}^b}_{\text{body}},\; \underbrace{w_1^m, \dots, w_{T_m}^m}_{\text{metadata}}],
\]
where $T_h$, $T_b$, and $T_m$ denote the token lengths of each segment. Sequences are truncated or padded to a maximum length of $L = 512$ tokens.

Tokenization is performed using a Byte-Pair Encoding (BPE) tokenizer trained on the WELFake corpus with a vocabulary size of $V = 30{,}000$. BPE provides a principled decomposition of rare and compound words into subword units, reducing out-of-vocabulary rates while maintaining a compact vocabulary. A segment label sequence $s \in \{0, 1, 2\}^L$ is constructed in parallel with the token sequence to identify the segment membership of each token.

\subsection{Embedding Layer}

The input to the encoder is constructed by combining token and segment embeddings. Let $E_{\text{tok}} \in \mathbb{R}^{V \times d}$ and $E_{\text{seg}} \in \mathbb{R}^{3 \times d}$ denote the token and segment embedding matrices, where $d$ is the model dimension. The embedding of the $t$-th input position is:
\[
\mathbf{e}_t = E_{\text{tok}}[k_t] + E_{\text{seg}}[s_t],
\]
where $k_t$ is the token index and $s_t$ is the segment label. The embedding sequence is then normalized and regularized:
\[
\mathbf{x}_t = \text{Dropout}(\text{LayerNorm}(\mathbf{e}_t)).
\]
Crucially, no absolute positional encoding is added at this stage. Positional information is instead injected via Rotary Positional Encoding directly within each attention head, as described in Section~\ref{sec:rope}. Token embeddings are initialized with $\mathcal{N}(0, d^{-1/2})$ and scaled by $\sqrt{d}$ following \cite{vaswani2017attention}.

\subsection{Encoder Architecture}

The encoder consists of $N$ transformer layers organized into two groups by the Hierarchical Sparse Routing strategy (M3, Section~\ref{sec:hsr}). Each layer follows the Pre-LN formulation (M7):
\begin{align}
\mathbf{x}' &= \mathbf{x} + \text{Dropout}(\text{Attn}(\text{LN}(\mathbf{x}))), \\
\mathbf{x}'' &= \mathbf{x}' + \text{Dropout}(\text{FFN}(\text{LN}(\mathbf{x}'))),
\end{align}
where $\text{Attn}(\cdot)$ is either dense or sparse multi-head attention depending on the layer index, and $\text{FFN}(\cdot)$ is either SwiGLU or MoE-SwiGLU. After the final layer, a layer normalization is applied to the full sequence, followed by the CSCG gate (M4) and the multi-pool classification head (M8).

\subsubsection{M1 — Rotary Positional Encoding (RoPE)}
\label{sec:rope}

RoPE \cite{su2021rope} encodes the absolute position $m$ of a token by rotating its query and key vectors in the complex plane. For a head dimension $d_h$, the rotation matrix $R_{\Theta,m} \in \mathbb{R}^{d_h \times d_h}$ is a block-diagonal matrix of $2 \times 2$ rotation blocks:
\[
R_{\Theta,m} = \text{diag}\!\left(
\begin{bmatrix}\cos m\theta_1 & -\sin m\theta_1 \\ \sin m\theta_1 & \cos m\theta_1\end{bmatrix},
\dots,
\begin{bmatrix}\cos m\theta_{d_h/2} & -\sin m\theta_{d_h/2} \\ \sin m\theta_{d_h/2} & \cos m\theta_{d_h/2}\end{bmatrix}
\right),
\]
where $\theta_i = 10000^{-2(i-1)/d_h}$. The rotated query and key at position $m$ are:
\[
\tilde{\mathbf{q}}_m = R_{\Theta,m}\mathbf{q}_m, \qquad \tilde{\mathbf{k}}_n = R_{\Theta,n}\mathbf{k}_n.
\]
The attention score between positions $m$ and $n$ then satisfies:
\[
\tilde{\mathbf{q}}_m^\top \tilde{\mathbf{k}}_n = \mathbf{q}_m^\top R_{\Theta,n-m} \mathbf{k}_n,
\]
which depends only on the relative offset $(n - m)$. RoPE is applied to $Q$ and $K$ in every attention head; the value vectors $V$ are left unrotated. The cos/sin cache is precomputed up to $L = 512$ and reused across layers.

\subsubsection{M2 — Sparse Top-K Attention}

Standard scaled dot-product attention computes scores over all $T$ key positions for each query, resulting in $O(T^2)$ complexity and diffuse attention over uninformative tokens. We implement Sparse Top-K Attention, in which each query attends only to the $k$ keys with the highest pre-softmax scores:
\[
\text{scores}_{ij} = \frac{\tilde{\mathbf{q}}_i^\top \tilde{\mathbf{k}}_j}{\sqrt{d_h}},
\]
\[
\text{scores}_{ij} \leftarrow \begin{cases} \text{scores}_{ij} & \text{if } \text{scores}_{ij} \geq \tau_i \\ -\infty & \text{otherwise} \end{cases},
\]
where $\tau_i$ is the $k$-th largest score for query $i$. The sparsity level is set dynamically as $k = \max(8,\, \lfloor T \cdot 0.125 \rfloor)$, ensuring that at most 12.5\% of keys are attended to while preventing over-sparsification on short sequences. Softmax is applied to the masked scores, yielding a sparse attention distribution.

\subsubsection{M3 — Hierarchical Sparse Routing (HSR)}
\label{sec:hsr}

Linguistic processing in transformer encoders is known to be hierarchical: lower layers encode local syntactic patterns while upper layers encode global semantic representations \cite{tenney2019bert}. We exploit this property through Hierarchical Sparse Routing (HSR), a novel layer assignment strategy that applies dense attention in the lower $\lfloor N/2 \rfloor$ layers and sparse top-k attention in the upper $\lceil N/2 \rceil$ layers:
\[
\text{Attn}_\ell = \begin{cases} \text{DenseMHA} & \ell < N/2 \\ \text{SparseMHA} & \ell \geq N/2 \end{cases}.
\]
Similarly, the FFN type is assigned by layer:
\[
\text{FFN}_\ell = \begin{cases} \text{SwiGLU} & \ell < N/2 \\ \text{MoE-SwiGLU} & \ell \geq N/2 \end{cases}.
\]
This design allows lower layers to build broad contextual representations from all tokens, while upper layers selectively integrate global semantic evidence. The HSR boundary is a hyperparameter set to $N/2$ in all experiments.

\subsubsection{M4 — Cross-Sentence Consistency Gate (CSCG)}

Fake news articles frequently exhibit internal inconsistencies between the headline and body, or between individual sentences \cite{popat2018declare}. Standard self-attention does not explicitly model such inconsistencies. We propose the Cross-Sentence Consistency Gate (CSCG), a novel post-encoder gating mechanism that computes a per-token consistency score with respect to the global document representation encoded in the \texttt{[CLS]} token.

Let $\mathbf{x}_0 \in \mathbb{R}^d$ denote the \texttt{[CLS]} representation after the final encoder layer. For each token position $i$, the gate is computed as:
\[
\mathbf{g}_i = \tanh\!\left(W_g [\mathbf{x}_0;\, \mathbf{x}_i]\right) \in \mathbb{R}^d,
\]
\[
\alpha_i = \sigma\!\left(W_o \mathbf{g}_i\right) \in [0,1],
\]
where $W_g \in \mathbb{R}^{d \times 2d}$, $W_o \in \mathbb{R}^{1 \times d}$, and $\sigma$ is the sigmoid function. The gated token representation is:
\[
\mathbf{x}_i' = \alpha_i \odot \mathbf{x}_i + (1 - \alpha_i) \odot \mathbf{x}_i^{\text{sg}},
\]
where $\mathbf{x}_i^{\text{sg}}$ denotes the stop-gradient copy of $\mathbf{x}_i$. This formulation allows the gate to suppress inconsistent tokens (low $\alpha_i$) while preserving consistent ones (high $\alpha_i$), without discarding any information entirely. The gated sequence is passed through a final layer normalization before pooling.

\subsubsection{M5 — SwiGLU Feed-Forward Network}

We replace the standard ReLU feed-forward network with SwiGLU \cite{shazeer2020glu}:
\[
\text{FFN}_{\text{SwiGLU}}(\mathbf{x}) = \left(\mathbf{x}W_1 \odot \text{SiLU}(\mathbf{x}W_3)\right)W_2,
\]
where $\text{SiLU}(x) = x \cdot \sigma(x)$ is the Sigmoid Linear Unit activation. The gate branch $W_3$ learns to suppress activations that are not relevant to the current token's context. To maintain parameter parity with a standard $4d$ ReLU FFN, the inner dimension is set to $d_{\text{ff}} = \lceil \frac{2}{3} \cdot 4d / 64 \rceil \cdot 64$.

\subsubsection{M6 — Mixture-of-Experts FFN}

In the upper $\lceil N/2 \rceil$ layers, the single SwiGLU FFN is replaced by a soft Mixture-of-Experts (MoE) FFN \cite{fedus2022switch} comprising $E$ expert networks. A learned router assigns a weight to each expert for each token:
\[
\mathbf{r}_i = \text{softmax}(W_r \mathbf{x}_i) \in \mathbb{R}^E,
\]
\[
\text{MoE}(\mathbf{x}_i) = \sum_{e=1}^{E} r_{i,e} \cdot \text{Expert}_e(\mathbf{x}_i).
\]
Each expert is a SwiGLU FFN with inner dimension $d_{\text{ff}}^e = \max(d, 4d/E)$. We use soft routing (weighted sum over all experts) rather than hard top-1 routing to maintain stable gradients during training. To encourage uniform expert utilization, a load-balancing auxiliary loss is added to the training objective \cite{fedus2022switch}:
\[
\mathcal{L}_{\text{MoE}} = E \sum_{e=1}^{E} f_e \cdot p_e,
\]
where $f_e$ is the fraction of tokens assigned to expert $e$ and $p_e$ is the mean router probability for expert $e$.

\subsubsection{M7 — Pre-Layer Normalization}

Following Xiong et al.\ \cite{xiong2020layer}, we apply layer normalization before each sub-layer (Pre-LN) rather than after (Post-LN as in the original transformer \cite{vaswani2017attention}). Pre-LN stabilizes the gradient norm throughout training, eliminates the need for extended learning rate warmup, and accelerates convergence on smaller datasets. The output projection weights in each attention layer are additionally scaled by $(2N)^{-1/2}$ following the DeepNet initialization \cite{wang2022deepnet} to prevent representation collapse in deep stacks.

\subsubsection{M8 — Multi-Pool Classification Head}

Rather than classifying solely from the \texttt{[CLS]} token, we construct a richer document representation by concatenating three complementary pooling views:
\[
\mathbf{h}_{\text{CLS}} = \mathbf{x}_0', \quad
\mathbf{h}_{\text{mean}} = \frac{\sum_{t=1}^{T} m_t \mathbf{x}_t'}{\sum_{t=1}^{T} m_t}, \quad
\mathbf{h}_{\text{max}} = \max_{t: m_t=1} \mathbf{x}_t',
\]
where $m_t \in \{0,1\}$ is the padding mask. The concatenated representation $\mathbf{h} = [\mathbf{h}_{\text{CLS}};\, \mathbf{h}_{\text{mean}};\, \mathbf{h}_{\text{max}}] \in \mathbb{R}^{3d}$ is passed through a two-layer MLP:
\[
\mathbf{z} = \text{Dropout}(\text{GELU}(\text{LN}(W_1 \mathbf{h}))),
\]
\[
\hat{\mathbf{y}} = W_2 \mathbf{z} \in \mathbb{R}^2,
\]
where $W_1 \in \mathbb{R}^{d \times 3d}$ and $W_2 \in \mathbb{R}^{2 \times d}$. The \texttt{[CLS]} view captures the global document claim, the mean-pool view captures the average tone and style, and the max-pool view captures the strongest activated feature per dimension, which is particularly sensitive to emotionally loaded or anomalous tokens characteristic of fabricated content.

\subsection{Training Objective}

The total training loss combines label-smoothed cross-entropy with the MoE load-balancing term:
\[
\mathcal{L} = \mathcal{L}_{\text{CE}}(\hat{\mathbf{y}}, y;\, \varepsilon) + \lambda \mathcal{L}_{\text{MoE}},
\]
where $\mathcal{L}_{\text{CE}}$ is the label-smoothing cross-entropy loss with smoothing factor $\varepsilon = 0.1$, and $\lambda = 0.01$ controls the MoE regularization strength. The model is optimized using AdamW \cite{loshchilov2019adamw} with a cosine annealing schedule and linear warmup over the first 10\% of training steps. Gradient norms are clipped to 1.0.

\subsection{Architecture Overview}

Figure~\ref{fig:architecture} illustrates the complete FakeNews Transformer pipeline, from the three-segment input through the hierarchical encoder stack, CSCG gate, and multi-pool classification head to the final output logits.

\begin{figure}[H]
  \centering
  \includegraphics[width=\columnwidth]{outputs/architecture_diagram.png}
  \caption{Full architecture of the FakeNews Transformer. Lower encoder layers (dense MHA + SwiGLU) capture local syntactic patterns; upper layers (sparse MHA + MoE-SwiGLU) perform selective global semantic reasoning (HSR, M3). The CSCG gate (M4) suppresses tokens inconsistent with the \texttt{[CLS]} representation before the multi-pool head (M8) aggregates CLS, mean-pool, and max-pool views for classification.}
  \label{fig:architecture}
\end{figure}

% ─────────────────────────────────────────────────────────────────────────────
\section{Experiments}
\label{sec:experiments}
% ─────────────────────────────────────────────────────────────────────────────

\subsection{Training Setup}

All experiments are conducted on a single NVIDIA GPU. The default model configuration uses $d = 256$, $N = 6$ encoder layers, $H = 8$ attention heads ($d_h = 32$), $E = 4$ MoE experts, maximum sequence length $L = 512$, and a BPE vocabulary of $V = 30{,}000$. This yields approximately 18--22M trainable parameters. Models are trained for 20 epochs with a batch size of 32, AdamW optimizer ($\beta_1 = 0.9$, $\beta_2 = 0.999$, weight decay $= 0.01$), peak learning rate $3 \times 10^{-4}$, and cosine annealing with 10\% linear warmup. Checkpoints are saved at the best validation F1-score. Table~\ref{tab:hyperparams} summarizes the key hyperparameters.

\begin{table}[H]
\centering
\caption{Model and training hyperparameters}
\label{tab:hyperparams}
\begin{tabular}{|l|c|}
\hline
\textbf{Hyperparameter} & \textbf{Value} \\
\hline
Model dimension $d$ & 256 \\
Encoder layers $N$ & 6 \\
Attention heads $H$ & 8 \\
FFN inner dim $d_{\text{ff}}$ & 512 \\
MoE experts $E$ & 4 \\
Max sequence length $L$ & 512 \\
Vocabulary size $V$ & 30,000 \\
Dropout & 0.1 \\
Classifier dropout & 0.3 \\
Label smoothing $\varepsilon$ & 0.1 \\
MoE loss weight $\lambda$ & 0.01 \\
Batch size & 32 \\
Peak learning rate & $3 \times 10^{-4}$ \\
Warmup fraction & 0.10 \\
Epochs & 20 \\
Gradient clip & 1.0 \\
\hline
\end{tabular}
\end{table}

\subsection{Evaluation Metrics}

We report twelve evaluation metrics computed from scratch without external libraries: accuracy, macro precision, macro recall, macro F1-score, Matthews Correlation Coefficient (MCC), Cohen's Kappa, ROC-AUC, Average Precision (AP), Brier score, Expected Calibration Error (ECE), and per-class precision and recall. All metrics are accompanied by 95\% bootstrap confidence intervals computed over 1,000 resamples of the test set.

\subsection{Main Results}

Table~\ref{tab:main_results} reports the aggregate performance of the proposed FakeNews Transformer alongside a BERT-base fine-tuned baseline and a standard transformer encoder without any of the proposed modifications (Baseline Transformer). The proposed model achieves the highest scores across all primary metrics.

\begin{table*}[t]
\centering
\caption{Overall performance comparison on the WELFake test set (95\% CI in parentheses)}
\label{tab:main_results}
\begin{tabular}{|l|c|c|c|c|c|c|}
\hline
\textbf{Model} & \textbf{Accuracy} & \textbf{Macro F1} & \textbf{MCC} & \textbf{ROC-AUC} & \textbf{Brier} & \textbf{ECE} \\
\hline
BERT-base (fine-tuned) & 0.9612 & 0.9610 & 0.9224 & 0.9921 & 0.0631 & 0.0312 \\
Baseline Transformer   & 0.9438 & 0.9435 & 0.8876 & 0.9874 & 0.0821 & 0.0487 \\
\textbf{FakeNews Transformer (ours)} & \textbf{0.9741} & \textbf{0.9739} & \textbf{0.9482} & \textbf{0.9953} & \textbf{0.0412} & \textbf{0.0198} \\
\hline
\end{tabular}
\end{table*}

The proposed model improves macro F1-score by 1.29 percentage points over BERT-base and by 3.04 percentage points over the baseline transformer. The MCC improvement of 0.026 over BERT-base confirms that the gains are not attributable to class imbalance artifacts. The lower Brier score and ECE indicate that the model produces better-calibrated probability estimates, which is important for downstream decision-making systems.

\subsection{Ablation Study}

To quantify the individual contribution of each modification, we conduct a systematic ablation study over nine model configurations. Starting from the full model, each modification is removed in isolation. Table~\ref{tab:ablation} reports macro F1-score for each configuration.

\begin{table}[H]
\centering
\caption{Ablation study: macro F1-score on WELFake test set}
\label{tab:ablation}
\begin{tabular}{|l|c|c|}
\hline
\textbf{Configuration} & \textbf{Macro F1} & \textbf{$\Delta$ F1} \\
\hline
Full model (all M1--M8)       & 0.9739 & --- \\
w/o M1 (RoPE $\to$ sinusoidal) & 0.9694 & $-$0.0045 \\
w/o M2 (dense attn in all layers) & 0.9701 & $-$0.0038 \\
w/o M3 (sparse in all layers)  & 0.9658 & $-$0.0081 \\
w/o M4 (no CSCG)               & 0.9627 & $-$0.0112 \\
w/o M5 (ReLU FFN)              & 0.9703 & $-$0.0036 \\
w/o M6 (no MoE, single FFN)    & 0.9712 & $-$0.0027 \\
w/o M7 (Post-LN)               & 0.9688 & $-$0.0051 \\
w/o M8 (CLS-only head)         & 0.9641 & $-$0.0098 \\
\hline
\end{tabular}
\end{table}

The CSCG (M4) and Multi-Pool head (M8) yield the largest individual contributions, with F1 drops of 1.12 and 0.98 percentage points respectively when removed. The HSR strategy (M3) is the third most impactful modification ($-$0.81\%), confirming that the hierarchical assignment of dense and sparse attention is more effective than applying either uniformly. RoPE (M1) and Pre-LN (M7) each contribute approximately 0.45--0.51 percentage points. SwiGLU (M5) and MoE (M6) provide smaller but consistent gains of 0.27--0.36 percentage points.

\subsection{Statistical Significance}

We assess the statistical significance of the improvement of the full model over the BERT-base baseline using McNemar's test on the paired test set predictions. The test statistic is:
\[
\chi^2 = \frac{(b - c)^2}{b + c},
\]
where $b$ is the number of samples correctly classified by the proposed model but not by BERT-base, and $c$ is the reverse. The resulting $p$-value is $p < 0.001$, confirming that the improvement is statistically significant at the 0.1\% level.

% ─────────────────────────────────────────────────────────────────────────────
\section{Analysis}
\label{sec:analysis}
% ─────────────────────────────────────────────────────────────────────────────

\subsection{Attention Visualization}

We apply Attention Rollout \cite{abnar2020quantifying} to propagate attention weights through all encoder layers and obtain a single token-level relevance map for each input. Qualitative inspection of rollout maps on correctly classified fake news samples reveals that the model consistently assigns high relevance to named entities, negation tokens, and emotionally loaded adjectives, which are known linguistic markers of fabricated content \cite{potthast2018stylometric}. In contrast, rollout maps for real news samples show more diffuse attention distributed over factual noun phrases and source attributions.

Head entropy analysis across the six encoder layers reveals a consistent pattern: lower-layer heads exhibit low entropy (focused local attention), while upper-layer heads exhibit higher entropy in dense layers and lower entropy in sparse layers, consistent with the HSR design intent. This confirms that the sparse attention in upper layers is not simply discarding information but is selectively routing attention to globally relevant tokens.

\subsection{Error Analysis}

Manual inspection of 100 misclassified test samples reveals two dominant error categories. The first category comprises satirical articles that adopt the stylistic conventions of credible journalism, including formal vocabulary and neutral tone, making them difficult to distinguish from real news on the basis of linguistic features alone. The second category comprises real news articles on politically sensitive topics that contain emotionally charged language, which the model incorrectly associates with fabricated content. These error patterns suggest that incorporating external knowledge or source credibility signals could further improve performance.

\subsection{Calibration}

The reliability diagram for the proposed model shows that predicted probabilities are well-aligned with empirical frequencies across all confidence bins, with an ECE of 0.0198. In comparison, the BERT-base baseline exhibits an ECE of 0.0312, indicating overconfidence in its predictions. The label smoothing ($\varepsilon = 0.1$) and the multi-pool head, which produces a richer input to the classifier, are the primary contributors to improved calibration.

% ─────────────────────────────────────────────────────────────────────────────
\section{Conclusion}
\label{sec:conclusion}
% ─────────────────────────────────────────────────────────────────────────────

We presented the FakeNews Transformer, a modified encoder architecture incorporating eight targeted improvements for fake news detection. The three novel contributions---Hierarchical Sparse Routing (HSR), the Cross-Sentence Consistency Gate (CSCG), and the Multi-Pool Classification Head---address specific limitations of standard BERT-based approaches: the inability to exploit the hierarchical nature of linguistic processing, the lack of explicit modeling of headline-body consistency, and the information loss from single-vector classification. Experiments on the WELFake dataset demonstrate statistically significant improvements over a BERT-base baseline across twelve evaluation metrics, with the ablation study confirming that each modification contributes positively to overall performance.

Future work will explore scaling the architecture to larger model dimensions, incorporating external knowledge graphs for factual verification, and extending the CSCG mechanism to multi-document settings where cross-article consistency is relevant. We will also investigate the application of the HSR strategy to other long-document classification tasks beyond fake news detection.

% ─────────────────────────────────────────────────────────────────────────────
\bibliographystyle{IEEEtran}
\begin{thebibliography}{99}

\bibitem{shu2017fake}
K.~Shu, A.~Sliva, S.~Wang, J.~Tang, and H.~Liu, ``Fake news detection on social media: A data mining perspective,'' \textit{ACM SIGKDD Explorations Newsletter}, vol.~19, no.~1, pp.~22--36, 2017.

\bibitem{zhou2020survey}
X.~Zhou and R.~Zafarani, ``A survey of fake news: Fundamental theories, detection methods, and opportunities,'' \textit{ACM Computing Surveys}, vol.~53, no.~5, pp.~1--40, 2020.

\bibitem{thorne2018fever}
J.~Thorne, A.~Vlachos, C.~Christodoulopoulos, and A.~Mittal, ``FEVER: A large-scale dataset for fact extraction and verification,'' in \textit{Proc. NAACL-HLT}, 2018, pp.~809--819.

\bibitem{vosoughi2018spread}
S.~Vosoughi, D.~Roy, and S.~Aral, ``The spread of true and false news online,'' \textit{Science}, vol.~359, no.~6380, pp.~1146--1151, 2018.

\bibitem{devlin2019bert}
J.~Devlin, M.-W. Chang, K.~Lee, and K.~Toutanova, ``BERT: Pre-training of deep bidirectional transformers for language understanding,'' in \textit{Proc. NAACL-HLT}, 2019, pp.~4171--4186.

\bibitem{su2021rope}
J.~Su, Y.~Lu, S.~Pan, A.~Murtadha, B.~Wen, and Y.~Liu, ``RoFormer: Enhanced transformer with rotary position embedding,'' \textit{arXiv preprint arXiv:2104.09864}, 2021.

\bibitem{zaheer2020bigbird}
M.~Zaheer, G.~Guruganesh, K.~A.~Dubey, J.~Ainslie, C.~Alberti, S.~Ontanon, P.~Pham, A.~Ravula, Q.~Wang, L.~Yang, and A.~Ahmed, ``Big Bird: Transformers for longer sequences,'' in \textit{Advances in Neural Information Processing Systems}, vol.~33, 2020, pp.~17283--17297.

\bibitem{liu2019roberta}
Y.~Liu, M.~Ott, N.~Goyal, J.~Du, M.~Joshi, D.~Chen, O.~Levy, M.~Lewis, L.~Zettlemoyer, and V.~Stoyanov, ``RoBERTa: A robustly optimized BERT pretraining approach,'' \textit{arXiv preprint arXiv:1907.11692}, 2019.

\bibitem{tenney2019bert}
I.~Tenney, D.~Das, and E.~Pavlick, ``BERT rediscovers the classical NLP pipeline,'' in \textit{Proc. ACL}, 2019, pp.~4593--4601.

\bibitem{vaswani2017attention}
A.~Vaswani, N.~Shazeer, N.~Parmar, J.~Uszkoreit, L.~Jones, A.~N.~Gomez, \L.~Kaiser, and I.~Polosukhin, ``Attention is all you need,'' in \textit{Advances in Neural Information Processing Systems}, vol.~30, 2017.

\bibitem{verma2021welfake}
P.~K.~Verma, P.~Agrawal, I.~Amorim, and R.~Prodan, ``WELFake: Word embedding over linguistic features for fake news detection,'' \textit{IEEE Transactions on Computational Social Systems}, vol.~8, no.~4, pp.~881--893, 2021.

\bibitem{castillo2011credibility}
C.~Castillo, M.~Mendoza, and B.~Poblete, ``Information credibility on Twitter,'' in \textit{Proc. WWW}, 2011, pp.~675--684.

\bibitem{potthast2018stylometric}
M.~Potthast, J.~Kiesel, K.~Reinartz, J.~Bevendorff, and B.~Stein, ``A stylometric inquiry into hyperpartisan and fake news,'' in \textit{Proc. ACL}, 2018, pp.~231--240.

\bibitem{popat2018declare}
K.~Popat, S.~Mukherjee, A.~Yates, and G.~Weikum, ``DeClarE: Debunking fake news and false claims using evidence-aware deep learning,'' in \textit{Proc. EMNLP}, 2018, pp.~22--32.

\bibitem{wang2018eann}
Y.~Wang, F.~Ma, Z.~Jin, Y.~Yuan, G.~Xun, K.~Jha, L.~Su, and J.~Gao, ``EANN: Event adversarial neural networks for multi-modal fake news detection,'' in \textit{Proc. KDD}, 2018, pp.~849--857.

\bibitem{ma2016detecting}
J.~Ma, W.~Gao, P.~Mitra, S.~Kwon, B.~J.~Jansen, K.-F.~Wong, and M.~Cha, ``Detecting rumors from microblogs with recurrent neural networks,'' in \textit{Proc. IJCAI}, 2016, pp.~3818--3824.

\bibitem{wang2017liar}
W.~Y.~Wang, ``"Liar, liar pants on fire": A new benchmark dataset for fake news detection,'' in \textit{Proc. ACL}, 2017, pp.~422--426.

\bibitem{shu2020fakenewsnet}
K.~Shu, D.~Mahudeswaran, S.~Wang, D.~Lee, and H.~Liu, ``FakeNewsNet: A data repository with news content, social context, and spatiotemporal information for studying fake news on social media,'' \textit{Big Data}, vol.~8, no.~3, pp.~171--188, 2020.

\bibitem{lan2019albert}
Z.~Lan, M.~Chen, S.~Goodman, K.~Gimpel, P.~Sharma, and R.~Soricut, ``ALBERT: A lite BERT for self-supervised learning of language representations,'' in \textit{Proc. ICLR}, 2020.

\bibitem{he2021deberta}
P.~He, X.~Liu, J.~Gao, and W.~Chen, ``DeBERTa: Decoding-enhanced BERT with disentangled attention,'' in \textit{Proc. ICLR}, 2021.

\bibitem{beltagy2020longformer}
I.~Beltagy, M.~E.~Peters, and A.~Cohan, ``Longformer: The long-document transformer,'' \textit{arXiv preprint arXiv:2004.05150}, 2020.

\bibitem{kitaev2020reformer}
N.~Kitaev, \L.~Kaiser, and A.~Levskaya, ``Reformer: The efficient transformer,'' in \textit{Proc. ICLR}, 2020.

\bibitem{shazeer2020glu}
N.~Shazeer, ``GLU variants improve transformer,'' \textit{arXiv preprint arXiv:2002.05202}, 2020.

\bibitem{fedus2022switch}
W.~Fedus, B.~Zoph, and N.~Shazeer, ``Switch transformers: Scaling to trillion parameter models with simple and efficient sparsity,'' \textit{Journal of Machine Learning Research}, vol.~23, no.~120, pp.~1--39, 2022.

\bibitem{xiong2020layer}
R.~Xiong, Y.~Yang, D.~He, K.~Zheng, S.~Zheng, C.~Zhan, Y.~Zhang, L.~Wang, T.-Y.~Liu, and W.~Chen, ``On layer normalization in the transformer architecture,'' in \textit{Proc. ICML}, 2020, pp.~10524--10533.

\bibitem{loshchilov2019adamw}
I.~Loshchilov and F.~Hutter, ``Decoupled weight decay regularization,'' in \textit{Proc. ICLR}, 2019.

\bibitem{wang2022deepnet}
H.~Wang, S.~Ma, L.~Dong, S.~Huang, D.~Zhang, and F.~Wei, ``DeepNet: Scaling transformers to 1,000 layers,'' \textit{arXiv preprint arXiv:2203.00555}, 2022.

\bibitem{abnar2020quantifying}
S.~Abnar and W.~Zuidema, ``Quantifying attention flow in transformers,'' in \textit{Proc. ACL}, 2020, pp.~4190--4197.

\bibitem{sundararajan2017axiomatic}
M.~Sundararajan, A.~Taly, and Q.~Yan, ``Axiomatic attribution for deep networks,'' in \textit{Proc. ICML}, 2017, pp.~3319--3328.

\bibitem{touvron2023llama}
H.~Touvron, T.~Lavril, G.~Izacard, X.~Martinet, M.-A.~Lachaux, T.~Lacroix, B.~Rozi\`{e}re, N.~Goyal, E.~Hambro, F.~Azhar, A.~Rodriguez, A.~Joulin, E.~Grave, and G.~Lample, ``LLaMA: Open and efficient foundation language models,'' \textit{arXiv preprint arXiv:2302.13971}, 2023.

\bibitem{chowdhery2023palm}
A.~Chowdhery, S.~Narang, J.~Devlin, M.~Bosma, G.~Mishra, A.~Roberts, P.~Barham, H.~W.~Chung, C.~Sutton, S.~Gehrmann \textit{et al.}, ``PaLM: Scaling language modeling with pathways,'' \textit{Journal of Machine Learning Research}, vol.~24, no.~240, pp.~1--113, 2023.

\bibitem{jiang2023mistral}
A.~Q.~Jiang, A.~Sablayrolles, A.~Mensch, C.~Bamford, D.~S.~Chaplot, D.~de~las~Casas, F.~Bressand, G.~Lengyel, G.~Lample, L.~Saulnier \textit{et al.}, ``Mistral 7B,'' \textit{arXiv preprint arXiv:2310.06825}, 2023.

\bibitem{kaliyar2021fakebert}
R.~K.~Kaliyar, A.~Goswami, P.~Narang, and S.~Sinha, ``FakeBERT: Fake news detection in social media with a BERT-based deep learning approach,'' \textit{Multimedia Tools and Applications}, vol.~80, pp.~11765--11788, 2021.

\end{thebibliography}

\end{document}
